%=============================================================================
% "Is there always time for everything?
%  How time allocation affects cognitive skill development during school years"
%
% Ricardo Antonio Avalos Uribe
% Economics paper format — converted from BA thesis (PUCP, 2021)
%=============================================================================

\documentclass[12pt]{article}

%--- Packages -----------------------------------------------------------------
\usepackage[margin=1in]{geometry}
\usepackage{amsmath,amssymb}
\usepackage{setspace}
\usepackage{booktabs}
\usepackage{array}
\usepackage{graphicx}
\usepackage[round,authoryear]{natbib}
\usepackage{hyperref}
\usepackage{microtype}
\usepackage{caption}

%--- Spacing ------------------------------------------------------------------
\onehalfspacing

%--- Notation convenience macros ---------------------------------------------
% NDC = Cognitive Development Level (Nivel de Desarrollo Cognitivo)
\newcommand{\NDC}{\mathrm{NDC}}
% Subscript shorthands
\newcommand{\child}{\text{child}}
\newcommand{\cg}{\text{cg}}   % caregiver

%=============================================================================
\begin{document}
%=============================================================================

%--- Title block --------------------------------------------------------------
\title{\textbf{Is there always time for everything?}\\
       \large How time allocation affects cognitive skill development\\
       during school years}

\author{Ricardo Antonio Avalos Uribe\thanks{
  Email: \href{mailto:ricardo.avalos@pucp.edu.pe}{ricardo.avalos@pucp.edu.pe}.
  I am grateful to Jos\'{e} Rodr\'{i}guez Gonz\'{a}les for guidance and
  to the Young Lives study for making the data available.}}

\date{December 2021}

\maketitle
\thispagestyle{empty}

%--- Abstract -----------------------------------------------------------------
\begin{abstract}
\noindent
This paper examines how children's time allocation across competing
activities—work, household chores, caregiving, and leisure—affects the
development of their cognitive skills during school years. Drawing on the
household decision-making framework and the educational production function,
I derive a value-added empirical specification and estimate it using five
rounds (2002–2016) of the Young Lives longitudinal survey for Peru.
The sample consists of adolescents aged approximately 15 in Round~5, and
cognitive development is measured by the Peabody Picture Vocabulary Test
(PPVT). Results show that time devoted to household chores and work outside
the home has a significant negative contemporaneous effect on PPVT scores.
Long-term negative effects are found for chores and early-age work (around
age~8). When the sample is disaggregated by sex and area of residence,
household chores primarily affect urban girls, while work primarily affects
boys—with contemporaneous effects for rural boys and long-term effects for
urban boys. These findings motivate low-cost informational policies targeting
parental awareness of the costs of early child labor and intensive household
task assignments.

\medskip
\noindent\textbf{JEL Codes:}~I24, J22, O15, D13

\noindent\textbf{Keywords:}~time use, cognitive development, child labor,
household chores, educational production function, Peru, value-added model
\end{abstract}

\newpage
\setcounter{page}{1}

%=============================================================================
\section{Introduction}
\label{sec:intro}
%=============================================================================

Education is arguably one of the most relevant topics when discussing a
nation's economic and social development. A particular view from the
perspective of human capital is that developed by \citet{becker1964}, who
suggests that education and the development of technical skills are themselves
human capital, since they allow access to future economic returns.
Subsequently, education as human capital would be extrapolated to its
importance in macroeconomic terms as a main component of growth and production
\citep{barro2002, romer1989}. On the other hand, a more holistic view suggests
that education or the development of cognitive skills is fundamental to social
development not only due to its economic returns, but also because it improves
the decision-making process regarding health, resource usage, and morality
\citep{busemeyer2018}. In short, education plays a fundamental role in
development in its broadest sense.

Related to these perspectives, Peru experienced significant progress in
education access: attendance rates in primary and secondary education in 2019
were around 97\% and 82\%, respectively. Additionally, school dropout rates
fell from 20\% to 6\% for primary school, and from 24\% to 11\% for secondary
school when comparing 2000 with 2015. This shows that access to education has
significantly improved. However, based on knowledge acquisition, the current
national status does not seem positive. Regarding expected learning outcomes
in reading comprehension, only 50\% of primary school students meet the
established expectation, while this figure drops to just 15\% for secondary
school students. The PISA tests present a similar scenario: compared to its
peers in the South American region, Peru ranks last in all learning
competencies \citep{minedu2019}. This reveals an urgent need to address the
educational challenge in terms of learning and skills acquisition, beyond
general education access rates.

Local literature has emphasized that living conditions, health status,
infrastructure, and teacher quality are key determinants of cognitive
development. On the one hand, poverty and poor nutrition during early childhood
have a negative effect on learning acquisition \citep{grantham2007,
sanchez2013}. On the other hand, studies related to the impact of school
infrastructure on learning outcomes point to moderate impacts, especially in
rural contexts where there is a lack of basic school infrastructure
\citep{campana2014}.

All of these research efforts contribute to a better understanding of learning
acquisition. However, factors more endogenous to the household and its
members—such as motivation, perceptions of the importance of education, and
time use—have not been as prominent in studies addressing cognitive
development. The use of time by school-aged children is particularly relevant
to this matter. Given that time is a scarce resource, when learning activities
are shared with other activities such as work, chores, or caregiving, cognitive
development could be jeopardized. This is relevant in Peru considering that
the school-grade repetition rate among children who work is 30\%, while among
non-working children it is 18\%. These statistics reveal a negative association
between children's time on activities outside of school and their learning
outcomes. Some qualitative studies indicate that this occurs especially during
adolescence, where the type of working activity is usually exposed to more
risks as the child starts joining the informal labor market
\citep{rojas2013}.

In this context, this research seeks to contribute to the national research
literature by providing empirical evidence on how children's time use affects
the development of their cognitive skills. Furthermore, it explores how
household geographical location and gender differences influence time use and,
consequently, impact cognitive development. The paper is structured as follows.
Section~\ref{sec:framework} presents a conceptual review emphasizing the
tradition of household decision-making models and the educational production
function. Section~\ref{sec:hypotheses} states the main hypotheses.
Section~\ref{sec:data} describes the data and stylized facts.
Section~\ref{sec:methodology} derives and discusses the econometric
specification. Section~\ref{sec:results} presents the estimation results.
Section~\ref{sec:conclusions} concludes with reflections and policy proposals.

%=============================================================================
\section{Conceptual Framework}
\label{sec:framework}
%=============================================================================

%-----------------------------------------------------------------------------
\subsection{Household decisions about time use}
\label{subsec:timeuse}
%-----------------------------------------------------------------------------

Time is a resource in the household that is considered scarce for all members.
Just as income has different uses, time can likewise be allocated to one purpose
or another depending on preferences and exogenous factors. A first approach to
these decisions and their connection to human capital formation was developed by
\citet{becker1994} in his work \textit{Human Capital}. Becker develops a
conceptualization where the use of time by each household member will have an
opportunity cost, which by convention would be the market wage. In this sense,
any use of time allocated for an activity other than work will have an
associated cost equal to the lost wage. This leads to the explicit inclusion
of time dedicated to education as human capital formation, where the associated
cost is offset by the higher wage expected in the future. This model allows
for the analysis of different contexts of the household and its members
regarding their use of time. For example, if a household has scarce economic
resources, the benefit of all its members contributing wages to family income
would have to be compared with the future benefit of some of them allocating
time to learning (human capital formation).

Within that framework, multiple empirical studies have attempted to test the
contexts in which children allocate time to work and education. It has been
observed that the time children allocate to work does not increase when wages
are high in the labor market because parents can find higher-paying jobs and
education would allow access to higher-paying jobs in the future
\citep{laszlo2008}. This is aligned with the Beckerian expectation extension,
where time decisions depend not only on current benefits but also on
perceptions of returns. Some observational studies have found that decisions
regarding schooling and time devoted to learning depend on perceptions of
returns, which in turn depend on the economic situation of the geographic
location and the labor market \citep{jensen2010, attanasio2009}. In addition,
qualitative evidence supports that in households where parents are self-employed
in agriculture, commerce, textiles, and other fields, there is a higher
prevalence of children's participation in family work \citep{rojas2013}. It
is therefore essential to understand that the basic determinants of the
Beckerian model—wages and future returns—are circumscribed by household
characteristics, and thus time-use decisions related to work and education
depend on these factors.

In addition to decisions about time spent on work and learning, there are also
decisions about how to use time for activities such as housework or caregiving.
Similar to time spent on education, based on the Beckerian model, decisions
about spending time on household chores and caregiving depend primarily on
potential income; however, they also depend on individual characteristics such
as sex, age, and household characteristics \citep{garcia2007}. In this regard,
empirical evidence indicates that female children spend most of their time on
household activities, and even when they do spend some time working, they spend
more time on household chores, with a higher prevalence among adolescents in
rural households \citep{garcia2007, zapata2010}. This shows that decisions
about time use, particularly regarding housework, are significantly determined
by gender roles and family characteristics.

%-----------------------------------------------------------------------------
\subsection{The generation of cognitive skills}
\label{subsec:cogskills}
%-----------------------------------------------------------------------------

The development of cognitive skills is understood as a process that depends on
a series of factors directly related to the environment through which learning
is acquired (the school and related aspects), as well as external factors. In
the international literature, this process is typically analyzed through the
educational production function. For example, \citet{glewwe2002}, building on
the Cobb--Douglas production function, specifies the inputs that determine
educational outcomes. He proposes three groups of fundamental factors: (i)
individual or household characteristics (demand); (ii) the quality of the
school and teachers (supply); and (iii) the individual's years of schooling.
This theoretical relationship allows us to describe what types of factors are
involved in the process of human capital formation.

This specification is well supported empirically. Many local and international
studies use this framework as a basis for estimation. On the one hand,
international evidence suggests that supply-side factors are of utmost
importance for educational performance. A meta-analysis by \citet{fuller1987}
shows that the quality of educational materials, expenditure per student, and
teacher quality are estimated to have a positive impact in the majority of
empirical studies (more than 50\%). On the other hand, local evidence suggests
that supply-side factors such as school equipment, type of institution, and
perceptions of teacher quality are highly important at all levels. Similarly,
regarding demand, there is some consensus on factors such as gender, native
language, and access to basic services \citep{miranda2008, jopen2014}. In
short, both supply and demand conditions must be considered when analyzing
educational performance, and the educational production function framework
allows us to elucidate which determinants to include in the analysis.

%-----------------------------------------------------------------------------
\subsection{Time distribution as a determinant of cognitive skill development}
\label{subsec:timedeterminant}
%-----------------------------------------------------------------------------

As previously reviewed, children's time allocation is directed toward different
activities that frequently compete for prevalence. Among these activities is
time dedicated to education or learning, which is a fundamental component for
the acquisition of cognitive skills. In this sense, any determinant that
directly affects time allocation will end up indirectly affecting the
development of cognitive skills.

In principle, under the models described in the previous sections, the use of
time in educational activities has a positive effect on cognitive development.
Meanwhile, time spent on other activities—such as household chores, family
work, or work outside the home—would have an effect that depends on the
household context or social determinants such as gender.

First, regarding the effect of time dedicated to education, evidence supports
the idea that more time devoted to learning allows for greater cognitive
development. Specifically, the more time devoted to learning in school or
outside of school, the lower the likelihood of grade repetition (a proxy for
learning acquisition). Additionally, this time dedicated by the child is
complemented by time dedicated by parents in the form of help or support in
the learning process \citep{beltran2013}. This shows that time dedicated to
studying—whether by the child himself or by the parents—has a positive effect
on learning outcomes. Consequently, any alternative use of time that limits
time in education would have an indirect and negative effect on cognitive
development.

Second, there is no consistent theoretical prediction regarding the effect of
hours spent working outside and inside the home, and empirical evidence on its
effects on cognitive development is mixed. In general, empirical results show
that longer hours spent working decrease cognitive development or learning
acquisition measured as academic delay \citep{beltran2013}. However, when
work is separated into that performed inside and outside the home, family work
begins to have a positive relationship with cognitive development, although
only in rural areas \citep{ponce2012}. Additionally, an improvement in wages
causes time spent working to decrease in favor of education in the presence
of greater future returns \citep{beltran2013}. Therefore, the type of work
activity, the structure of the family economy, and market conditions determine
the effect that hours spent working can have on cognitive development outcomes.

Finally, regarding hours dedicated to household activities, similar to any
other activity alternative to education, the effect on cognitive development
will depend on the intensity of the time dedicated as well as specific social
determinants. In this regard, the time dedicated to household activities is
greater for women than for men, both due to social structures and differences
in returns to the labor market \citep{garcia2007}. Additionally, the presence
of the mother at home or of a female figure in the home allows the time
dedicated to household chores to decrease, especially for girls, generating a
positive effect on schooling as they have more time for education
\citep{njie2015, levison1998}. Thus, regarding household chores, the
fundamental determinant of time distribution and, consequently, cognitive
development would be the family composition and the sex of the child.

In short, the conceptual and empirical evidence on time use and cognitive
development clearly demonstrates the existence of determinants beyond salary
as a counterpart to time use. On the one hand, the relationship between time
use in work activities (inside or outside the home) and the development of
cognitive skills will depend on the structure of the family economy and labor
market conditions. On the other hand, the relationship between time use in
household activities (caregiving and chores) and cognitive development is
closely related to family structure and gender roles.

%=============================================================================
\section{Hypothesis Statement}
\label{sec:hypotheses}
%=============================================================================

%-----------------------------------------------------------------------------
\subsection{General hypothesis}
\label{subsec:genhyp}
%-----------------------------------------------------------------------------

\textbf{General hypothesis:} The distribution of time among boys and girls
affects the development of their cognitive skills.

The development of cognitive skills, from the most general perspective through
the cognitive skills production function, is determined by a series of factors.
One of these factors is the allocation of time use, both because it determines
the time resources directly devoted to learning activities and because the use
of time in activities such as work or household chores has indirect impacts on
the development of cognitive skills.

%-----------------------------------------------------------------------------
\subsection{Hypothesis 1: Contemporary and long-term effects}
\label{subsec:hyp1}
%-----------------------------------------------------------------------------

\textbf{Hypothesis 1:} There are contemporaneous and long-term effects of
time use on the development of cognitive skills.

Evidence suggests that the distribution of children's time use affects the
development of their cognitive skills contemporaneously. Whether working inside
or outside the home, or performing domestic activities, the use of time in
activities other than learning has an immediate effect on the acquisition of
cognitive skills by limiting the time devoted to learning and by having
indirect effects on motivation or the perception of the benefits of education.
Furthermore, these contemporaneous effects have an impact on future cognitive
development, as reflected in the higher rates of school failure among those
who devote time to work or activities other than learning.

%-----------------------------------------------------------------------------
\subsection{Hypothesis 2: Differential effects by gender and area of residence}
\label{subsec:hyp2}
%-----------------------------------------------------------------------------

\textbf{Hypothesis 2:} There are differences by gender and area of residence
in the impact that time use has on cognitive development.

The impacts of time use on cognitive development are differentiated by the
child's gender and area of residence. This is particularly true because the
presence of gender roles makes home-based activities more prevalent for girls
and adolescents, which means this group sees their time devoted to education
limited and, therefore, achieves less cognitive development than their male
peers. Furthermore, the area of residence is crucial as it helps delineate
households where family-based economic structures predominate and where labor
markets are not as dynamic as to reduce the hours invested in education due
to lower returns.

%=============================================================================
\section{Data and Stylized Facts}
\label{sec:data}
%=============================================================================

%-----------------------------------------------------------------------------
\subsection{Sources of information}
\label{subsec:datasources}
%-----------------------------------------------------------------------------

The data source used is the international Young Lives study, known in Peru as
\textit{Ni\~{n}os del Milenio} (Children of the Millennium). This is an
international longitudinal study of child poverty that began in 2002 and
followed 12,000 children of low and middle income in four countries: Ethiopia,
India, Peru, and Vietnam. It was designed to understand how child poverty
impacts development. To achieve this, two age cohorts are followed throughout
childhood, adolescence, and adulthood, allowing comparisons of the same
individuals over time.

The study conducted five survey rounds between 2002 and 2016. During this
period, large-scale household surveys were conducted targeting all children
and their primary caregiver, in addition to in-depth interviews, group work,
and case studies with subsamples of children, parents or caregivers, teachers,
and community officials.

Since the emphasis of this paper is on time use and the development of
cognitive skills, some regularities deserve attention. On the one hand,
reports on time use come from two sources: primary caregivers (parents or
guardians) and the child themselves. For children under 12 years of age,
only the caregiver's report is available; those between 12 and 15 years of
age have both reports; and those over 15 years of age self-report their time
use. On the other hand, regarding information on cognitive development, there
are three types of tests: the Peabody Picture Vocabulary Test (PPVT), the
verbal test, and the mathematics test. Of these, only the PPVT is available
from the second round onwards (from age~5). The structure of the Young Lives
dataset is summarized in Appendix~\ref{app:data}.

%-----------------------------------------------------------------------------
\subsection{Sample and main variables}
\label{subsec:sample}
%-----------------------------------------------------------------------------

Given the study's objective, the youngest cohort from Round~5 (around 15 years
of age) was chosen as the main sample. This choice was based on the timeliness
of the data (2016) and because during this stage, time distribution is more
variable across different activity categories, as during adolescence there is
a tendency to engage in activities other than education, especially in adverse
contexts. It should be noted that while Round~5 is the main sample for
empirical results, variables from previous rounds are used to estimate
long-term effects.

\begin{table}[ht]
\centering
\caption{Main variables of the study}
\label{tab:mainvars}
\begin{tabular}{lp{9cm}}
\toprule
\textbf{Variable} & \textbf{Description} \\
\midrule
\multicolumn{2}{l}{\textit{Dependent variable}} \\
PPVT score & Score obtained in the Peabody Picture Vocabulary Test \\
\addlinespace
\multicolumn{2}{l}{\textit{Main independent variables}} \\
Child's leisure time       & Time spent on leisure activities (sleeping, eating, playing, etc.) \\
Child's time on chores     & Time spent on household chores (cleaning, tidying, cooking, etc.) \\
Child's time on caregiving & Time spent on caregiving activities for a family member at home \\
Child's time on work       & Time spent on work activities both inside and outside the family environment \\
\bottomrule
\end{tabular}
\end{table}

%-----------------------------------------------------------------------------
\subsection{Stylized facts}
\label{subsec:stylizedfacts}
%-----------------------------------------------------------------------------

To explore the feasibility of the proposed hypotheses, a descriptive analysis
of the relationship between cognitive development (measured by the PPVT score)
and the various time uses was conducted. Table~\ref{tab:descriptive} presents
the average PPVT score obtained in Round~5 by activity performed and the
rounds in which it was performed.

\begin{table}[ht]
\centering
\caption{Average PPVT score in Round 5 by time use and survey round of
activity}
\label{tab:descriptive}
\begin{tabular}{lcccccc}
\toprule
 & \multicolumn{6}{c}{\textbf{Round(s) in which the activity was performed}} \\
\cmidrule(lr){2-7}
\textbf{Activity} & None & R5 only & R4--5 & R3 only & R3--4 & R3--4--5 \\
\midrule
Chores
  & 87.01 & 80.58 & 80.75 & 81.68 & 80.20 & 74.27 \\
  & [29]  & [79]  & [335] & [21]  & [123] & [989] \\
\addlinespace
Caregiving
  & 78.87 & 79.34 & 76.06 & 76.69 & 74.44 & 73.20 \\
  & [600] & [134] & [244] & [89]  & [183] & [215] \\
\addlinespace
Family work
  & 80.26 & 77.12 & 72.98 & 71.85 & 67.00 & 63.61 \\
  & [1052]& [94]  & [91]  & [106] & [87]  & [53]  \\
\addlinespace
Market work
  & 77.38 & 69.84 & 74.13 & 73.87 &  ---  &  ---  \\
  & [1655]& [53]  & [3]   & [6]   & [0]   & [0]   \\
\bottomrule
\multicolumn{7}{l}{\footnotesize Notes: Number of observations in brackets.
Source: Young Lives. Authors' calculations.}
\end{tabular}
\end{table}

For all activities except market work, performing the activity contemporaneously
(in Round~5) is associated with a decrease in the PPVT score obtained in the
same round. This decrease is sustained if the activity has been performed since
Round~4, and a significant decrease occurs when the activity is performed
predominantly (from Round~3 to Round~5). Furthermore, a long-term descriptive
relationship is observed, since performing the activity only in Round~3 (age~8)
is associated with obtaining a lower PPVT score in Round~5 (age~15). Therefore,
this descriptive exercise shows evidence that the use of time in activities
other than education could have effects on the development of cognitive skills.
The objective of the empirical methodology is to obtain causal results to
confirm these preliminary observations.

%=============================================================================
\section{Methodology and Empirical Specification}
\label{sec:methodology}
%=============================================================================

%-----------------------------------------------------------------------------
\subsection{Derivation of the cognitive development equation}
\label{subsec:derivation}
%-----------------------------------------------------------------------------

To consistently analyze how time use affects the development of cognitive
skills, I start from the household's decision-making framework regarding time
use throughout the life cycle. This would be a two-member household—caregiver
and child—whose consumption depends not only on market goods $X_{t}^{j}$, but
also on the use of time to make such consumption effective
$T_{t}^{j}$ \citep{becker1993}. Furthermore, the child's level of cognitive
development $\NDC_{\child,t}$ is a component that gives utility to the
household and follows a process of accumulation throughout the life cycle
\citep{benporath1967}. The household optimization problem is:

\begin{equation}
  \max \; U = u(C_1, C_2, \ldots, C_T)
  \label{eq:utility}
\end{equation}

\begin{equation}
  C_t = f_c\!\left(\NDC_{\child,t},\; C_t^{TD},\; C_t^{O}\right)
  \label{eq:consumption}
\end{equation}

\begin{equation}
  C_t^{j} = f_j\!\left(X_t^{j},\, T_t^{j}\right)
  \qquad \forall\, j \in \{O,\, TD\}
  \label{eq:subconsumption}
\end{equation}

where $O$ denotes leisure and $TD$ denotes domestic tasks. The child's
cognitive development is governed by an accumulation equation:

\begin{equation}
  \NDC_{\child,t} = \NDC_{\child,t-1} + F(A_0,\, E_t,\, Q,\, N,\, H,\, I)
  \label{eq:ndc}
\end{equation}

\begin{equation}
  E_t = f_E\!\left(A_0,\; T_t^{E},\; X_t^{E},\; Q,\; N,\; H,\; P_t\right)
  \label{eq:education}
\end{equation}

where $A_0$ is innate ability, $E_t$ is educational achievement at period $t$,
$Q$ is school quality, $N$ are child characteristics, $H$ are household
characteristics, $I$ is an indicator of the informational environment, and
$P_t$ is the price level. The system is completed with the household budget
constraint:

\begin{equation}
  P_t X_t = Y_0 + Y\!\left(\hat{W}_t,\, T_t^{TM}\right)
            + Y\!\left(X_t^{TF},\, T_t^{TF},\, P_t\right)
  \label{eq:budget}
\end{equation}

where $TM$ denotes market work, $TF$ denotes family work, $\hat{W}_t$ is the
market wage, and $Y_0$ is autonomous income. The time constraint is:

\begin{equation}
  T = T_t^{O} + T_t^{TD} + T_t^{E} + T_t^{TF} + T_t^{TM}
  \label{eq:timeconstraint}
\end{equation}

A detailed derivation of the full system and its solution is presented in
Appendix~\ref{app:model}. The optimal time allocation by household members
depends on a set of exogenous factors. With this result, equation~\eqref{eq:ndc}
is evaluated at the optimum, yielding:

\begin{equation}
  \NDC_{\child,t} = \NDC_{\child,t-1}
    + F\!\left(T_t^{O(*)},\; T_t^{TD(*)},\; T_t^{TF(*)},\; T_t^{TM(*)},\;
               X_t^{E(*)},\; A_0,\; Q,\; N,\; H,\; I,\; P_t\right)
  \label{eq:ndcoptimal}
\end{equation}

Equation~\eqref{eq:ndcoptimal} provides a specification of the relationship
between the various uses of household members' time and their relationship to
the level of cognitive development. This is consistent with the household's
decision-making framework, since time is a limited resource and there will be
competition for time allocation. Consequently, more time devoted to leisure,
housework, or forms of work decreases the time actually devoted to school and,
consequently, to the development of cognitive skills.

%-----------------------------------------------------------------------------
\subsection{Empirical specification}
\label{subsec:empirical}
%-----------------------------------------------------------------------------

In order to empirically analyze the relationship between the level of cognitive
development and alternative uses of time, the function derived in
equation~\eqref{eq:ndcoptimal} is assumed to be linear. Furthermore,
considering that cognitive development is an accumulation process, a
specification can be derived that depends on the history of time distribution
and exogenous components, in addition to measurement error:

\begin{equation}
  \NDC_{\child,T} = \beta_0 + \beta_1\,\NDC_{\child,0}
    + \sum_{t=1}^{T} \sum_{j \in \{O,\,TD,\,TF,\,TM\}}
      T_{t}^{j(*)} \cdot \beta_{t}^{j}
    + \mathrm{Exogenous}_{t} \cdot \delta_{t}
    + \varepsilon_{T}
  \label{eq:valueadded}
\end{equation}

One possible method of direct estimation is ordinary least squares (OLS).
However, this would be biased due to omitted variables and selection. A direct
source of omitted-variable bias is the innate ability component $A_0$ and
household preferences for education. Both variables are unobservable and
directly impact the development of cognitive skills; their omission implies
biased estimates. On the other hand, both variables, being exogenous,
determine time allocation, generating selection bias since they are unobserved.
This concern is well documented in the literature of empirical estimates of
educational outcomes \citep{glewwe2020, webbink2005}.

Consequently, two viable estimation options are considered: a value-added type
estimation and estimation by instrumental variables. The value-added model is
proposed as the baseline option, since the information provided by the Young
Lives study allows for the adequate incorporation of a rich set of controls to
limit estimation bias. Furthermore, the specification in
equation~\eqref{eq:valueadded} directly accommodates the assumption that the
current outcome is determined by the accumulation of results from past periods.
Table~\ref{tab:modelvars} shows the variables used to build the model.

\begin{table}[ht]
\centering
\caption{Model variable specification}
\label{tab:modelvars}
\begin{tabular}{lp{9cm}}
\toprule
\textbf{Variable} & \textbf{Description} \\
\midrule
\multicolumn{2}{l}{\textit{Dependent variable}} \\
PPVT score & Score obtained in the Peabody Picture Vocabulary Test \\
\addlinespace
\multicolumn{2}{l}{\textit{Main independent variables}} \\
Child's leisure time       & Time spent on leisure (sleeping, eating, playing, etc.) \\
Child's time on chores     & Time spent on household chores (cleaning, tidying, cooking, etc.) \\
Child's time on caregiving & Time spent on caregiving activities for a family member \\
Child's time on work       & Time spent on work activities inside and outside the family \\
\addlinespace
\multicolumn{2}{l}{\textit{Additional controls}} \\
Child characteristics      & Nutritional status, preschool attendance, ongoing illness, etc. \\
Parental characteristics   & Age, highest level of education, legal marital status, etc. \\
School characteristics     & Teacher absence, travel time to school, type of institution, etc. \\
Household characteristics  & Region or area of residence, access to water and sanitation, wealth index, etc. \\
\bottomrule
\end{tabular}
\end{table}

%=============================================================================
\section{Estimation Results}
\label{sec:results}
%=============================================================================

%-----------------------------------------------------------------------------
\subsection{General results}
\label{subsec:general}
%-----------------------------------------------------------------------------

As a first result, the model is estimated for all individuals in Round~5
without subdivision by urban/rural area or by sex. From this perspective, the
estimated coefficients indicate the average relationship between time use and
the PPVT test result in Round~5. Table~\ref{tab:general} shows the results of
the standardized coefficients from this first estimation.

\begin{table}[ht]
\centering
\caption{General level estimate}
\label{tab:general}
\begin{tabular}{lc}
\toprule
 & \textbf{PPVT Score (Round 5)} \\
\midrule
\multicolumn{2}{l}{\textit{Round 5 (contemporaneous)}} \\
\quad Chores  & $-0.075^{**}$ \\
              & $(0.484)$ \\
\quad Caregiving & $-0.002$ \\
              & $(0.423)$ \\
\quad Leisure & $-0.086^{***}$ \\
              & $(0.256)$ \\
\quad Work    & $-0.043$ \\
              & $(0.359)$ \\
\addlinespace
\multicolumn{2}{l}{\textit{Round 4}} \\
\quad Chores  & $-0.003$ \\
              & $(0.585)$ \\
\quad Caregiving & $-0.005$ \\
              & $(0.444)$ \\
\quad Leisure & $-0.030$ \\
              & $(0.266)$ \\
\quad Work    & $-0.018$ \\
              & $(0.401)$ \\
\addlinespace
\multicolumn{2}{l}{\textit{Round 3 (long-term)}} \\
\quad Chores  & $-0.054^{*}$ \\
              & $(0.623)$ \\
\quad Caregiving & $-0.011$ \\
              & $(0.508)$ \\
\quad Leisure & $\phantom{-}0.007$ \\
              & $(0.280)$ \\
\quad Work    & $-0.061^{*}$ \\
              & $(0.648)$ \\
\addlinespace
$N$           & $1{,}188$ \\
$R^2$         & $0.455$ \\
Controls      & Yes \\
\bottomrule
\multicolumn{2}{l}{\footnotesize Standardized coefficients reported; standard
errors in parentheses.} \\
\multicolumn{2}{l}{\footnotesize $^{*}p<0.05$, $^{**}p<0.01$,
$^{***}p<0.001$.}
\end{tabular}
\end{table}

This first result highlights a long-term relationship between time use at work
and chores, and a contemporaneous relationship between time use on chores.
Specifically, the three aforementioned activities have a negative impact on
the measure of cognitive development. Regarding time use at work, the result
is consistent with the reviewed literature, as working at an early age (8
years) does not provide sufficient household income to compensate for the time
spent on such activities. On the other hand, chores also have a negative
impact both in the long term and contemporaneously. This makes sense because
such activities do not have a compensation—such as wages—to offset the effect
of the decrease in available time. Additionally, if the activity is performed
at an early and critical age (around 8 years), there will be limited time for
the development of cognitive skills.

%-----------------------------------------------------------------------------
\subsection{Results by area of residence and sex}
\label{subsec:heterogeneity}
%-----------------------------------------------------------------------------

A second set of results is presented in terms of different estimates by area
of residence and individual sex. This is because the literature suggests that
time allocation and educational outcomes diverge when these two variables are
combined. Table~\ref{tab:heterogeneity} presents the estimation results.

\begin{table}[ht]
\centering
\caption{Estimates by area of residence and sex}
\label{tab:heterogeneity}
\begin{tabular}{lcccc}
\toprule
 & \multicolumn{2}{c}{\textbf{Urban}} & \multicolumn{2}{c}{\textbf{Rural}} \\
\cmidrule(lr){2-3} \cmidrule(lr){4-5}
 & \textbf{Boy} & \textbf{Girl} & \textbf{Boy} & \textbf{Girl} \\
\midrule
\multicolumn{5}{l}{\textit{Round 5 (contemporaneous)}} \\
\quad Chores
  & $-0.007$       & $-0.109^{*}$  & $-0.120$      & $-0.127$ \\
  & $(0.863)$      & $(0.759)$     & $(2.080)$     & $(1.476)$ \\
\quad Caregiving
  & $\phantom{-}0.015$ & $\phantom{-}0.013$ & $-0.067$ & $-0.021$ \\
  & $(0.659)$      & $(0.696)$     & $(1.871)$     & $(1.433)$ \\
\quad Leisure
  & $-0.042$       & $-0.114^{*}$  & $\phantom{-}0.008$ & $-0.165$ \\
  & $(0.432)$      & $(0.424)$     & $(0.943)$     & $(0.902)$ \\
\quad Work
  & $-0.036$       & $-0.014$      & $-0.237^{**}$ & $\phantom{-}0.104$ \\
  & $(0.683)$      & $(0.618)$     & $(0.935)$     & $(1.251)$ \\
\addlinespace
\multicolumn{5}{l}{\textit{Round 4}} \\
\quad Chores
  & $\phantom{-}0.021$ & $\phantom{-}0.000$ & $-0.053$ & $\phantom{-}0.050$ \\
  & $(0.921)$      & $(0.987)$     & $(2.019)$     & $(1.987)$ \\
\quad Caregiving
  & $-0.023$       & $\phantom{-}0.013$ & $-0.005$ & $\phantom{-}0.032$ \\
  & $(0.750)$      & $(0.680)$     & $(1.813)$     & $(1.540)$ \\
\quad Leisure
  & $-0.059$       & $-0.029$      & $\phantom{-}0.006$ & $\phantom{-}0.023$ \\
  & $(0.409)$      & $(0.434)$     & $(1.035)$     & $(0.974)$ \\
\quad Work
  & $\phantom{-}0.045$ & $-0.012$  & $-0.136$      & $\phantom{-}0.092$ \\
  & $(0.730)$      & $(0.769)$     & $(1.044)$     & $(1.217)$ \\
\addlinespace
\multicolumn{5}{l}{\textit{Round 3 (long-term)}} \\
\quad Chores
  & $-0.065$       & $-0.027$      & $-0.229^{*}$  & $\phantom{-}0.029$ \\
  & $(0.997)$      & $(1.006)$     & $(2.211)$     & $(2.249)$ \\
\quad Caregiving
  & $\phantom{-}0.018$ & $-0.004$  & $\phantom{-}0.040$ & $-0.115$ \\
  & $(0.830)$      & $(0.834)$     & $(2.068)$     & $(1.666)$ \\
\quad Leisure
  & $-0.021$       & $-0.006$      & $-0.059$      & $\phantom{-}0.073$ \\
  & $(0.448)$      & $(0.423)$     & $(1.280)$     & $(0.998)$ \\
\quad Work
  & $-0.094^{*}$   & $-0.081$      & $-0.084$      & $\phantom{-}0.103$ \\
  & $(1.314)$      & $(1.583)$     & $(1.573)$     & $(1.581)$ \\
\addlinespace
$N$  & $472$  & $418$  & $142$  & $156$ \\
$R^2$ & $0.403$ & $0.449$ & $0.653$ & $0.628$ \\
Controls & Yes & Yes & Yes & Yes \\
\bottomrule
\multicolumn{5}{l}{\footnotesize Standardized coefficients reported; standard
errors in parentheses.} \\
\multicolumn{5}{l}{\footnotesize $^{*}p<0.05$, $^{**}p<0.01$,
$^{***}p<0.001$.}
\end{tabular}
\end{table}

For boys, the use of time spent working always has a negative effect.
Particularly in urban areas, this impact is only observed as a long-term
effect, consistent with the fact that early employment does not generate a
significant income counterpart and the type of work in urban areas tends to
limit cognitive development outcomes. In rural areas, the impact is rather
contemporaneous and is linked to the presence of more time- and effort-intensive
work, a recurring theme given the type of activity in rural areas. It is also
observed that chores have a negative effect in rural areas, probably because
they limit available time without an offsetting income or cognitive skills
improvement.

Among girls, statistically significant results are only observed in urban
settings and contemporaneously. The use of time on chores negatively impacts
PPVT test results. This could be related to gender archetypes, as girls around
15 years of age are charged with a greater burden of time and effort for such
activities, limiting their cognitive development.

In summary, the results disaggregated by sex and area of residence show that
the impact of time use is differentiated by these same variables. Girls are
particularly affected by household-related activities; while boys are affected
by work-related activities in the labor market or family environment, although
in rural areas household activities also affect them to some extent.
Consequently, the hypothesis of differentiated effects by sex and area of
residence is supported by the estimates shown.

%=============================================================================
\section{Conclusions and Policy Proposals}
\label{sec:conclusions}
%=============================================================================

This research analyzed how the allocation of time among children around 15
years of age affects their cognitive development. Based on previously
established evidence, it is suggested that the use of time in activities other
than educational activities affects educational outcomes and cognitive
development. Time- and effort-intensive activities such as work and housework
are particularly likely to have such an effect.

Two hypotheses were proposed and tested. On the one hand, time use has
long-term impacts on the development of cognitive skills. On the other hand,
time use has differential impacts depending on the area of residence (rural or
urban) and the individual's sex. Using the Young Lives longitudinal survey
and a value-added model, it is confirmed that there is indeed a negative
relationship between time use and cognitive development. Specifically, there
is a negative and contemporaneous relationship between time use on chores and
work (inside and outside the home) with respect to the cognitive development
measure used. Furthermore, when the sample is subdivided by sex and area of
residence, it is found that work activities affect boys, while chores
primarily affect girls.

These results show that cognitive development outcomes can be affected when
other activities such as chores and work are performed. Specifically, there is
a long-term relationship that holds true even when the sample is subdivided,
indicating that exposing boys and girls around the age of 8 to both activities
could have negative long-term consequences. In this sense, a possible
cost-effective policy would be to implement informational programs on the
benefits of education so that parents' preferences are focused on their
children's educational and cognitive development. Furthermore, information
could be included regarding the negative impacts generated when work or chores
are performed intensively at a young age (around 8 years old). In short, these
would be low-cost policies amenable to rigorous impact evaluation.

%=============================================================================
%  REFERENCES
%=============================================================================
\newpage

\begin{thebibliography}{99}

\bibitem[Attanasio and Kaufmann(2009)]{attanasio2009}
Attanasio, O., and K.~Kaufmann (2009).
\newblock Educational choices, subjective expectations, and credit constraints.
\newblock \textit{NBER Working Paper}.

\bibitem[Barro(2002)]{barro2002}
Barro, R. (2002).
\newblock Education as a determinant of economic growth.
\newblock \textit{Hoover Institution Press}.

\bibitem[Becker(1964)]{becker1964}
Becker, G. (1964).
\newblock Human capital revisited.
\newblock In G.~Becker, \textit{Human Capital}, pp.~15--21. Chicago:
  University of Chicago Press.

\bibitem[Becker(1993)]{becker1993}
Becker, G. (1993).
\newblock Investment in human capital: Rates of return.
\newblock In G.~Becker, \textit{A Theoretical and Empirical Analysis with
  Special Reference to Education}, pp.~59--92. Chicago: University of Chicago
  Press.

\bibitem[Becker(1994)]{becker1994}
Becker, G. (1994).
\newblock \textit{Human Capital: A Theoretical and Empirical Analysis}.
\newblock Chicago: University of Chicago Press.

\bibitem[Beltr\'{a}n(2013)]{beltran2013}
Beltr\'{a}n, A. (2013).
\newblock Family time is a scarce resource.
\newblock \textit{CIUP}, pp.~117--156.

\bibitem[Ben-Porath(1967)]{benporath1967}
Ben-Porath, Y. (1967).
\newblock The production of human capital and the life cycle of earnings.
\newblock \textit{Journal of Political Economy}, 352--362.

\bibitem[Busemeyer et~al.(2018)]{busemeyer2018}
Busemeyer, M., N.~Cloete, G.~Drori, L.~Lassnigg, and B.~Schober (2018).
\newblock The contribution of education to social progress.
\newblock \textit{Cambridge University Press}, pp.~753--778.

\bibitem[Campana et~al.(2014)]{campana2014}
Campana, Y., J.~Aguirre, D.~Velasco, and E.~Guerrero (2014).
\newblock Investment in educational infrastructure: The experience of flagship
  schools.
\newblock \textit{CIES}, pp.~6--13.

\bibitem[Fuller(1987)]{fuller1987}
Fuller, B. (1987).
\newblock What school factors raise achievement in the third world?
\newblock \textit{Review of Educational Research}, 255--292.

\bibitem[Garc\'{i}a(2007)]{garcia2007}
Garc\'{i}a, L. (2007).
\newblock Who does the chores? Estimation of a household production function
  in Peru.
\newblock \textit{MPRA}.

\bibitem[Glewwe(2002)]{glewwe2002}
Glewwe, P. (2002).
\newblock Schools and skills in developing countries: Education policies and
  socioeconomic outcomes.
\newblock \textit{Journal of Economic Literature}, 436--482.

\bibitem[Glewwe and Kremer(2006)]{glewwe2006}
Glewwe, P., and M.~Kremer (2006).
\newblock Schools, teachers and education outcomes in developing countries.
\newblock In E.~Hanushek and F.~Welch (Eds.), \textit{Handbook of the
  Economics of Education}, pp.~964--971.

\bibitem[Glewwe et~al.(2020)]{glewwe2020}
Glewwe, P., S.~Lambert, and C.~Chen (2020).

\bibitem[Gronau(1977)]{gronau1977}
Gronau, R. (1977).
\newblock Leisure, home production, and work---the theory of the allocation of
  time revisited.
\newblock \textit{Journal of Political Economy}, 85(6), 1099--1123.
\newblock Causal inference in education research.
\newblock \textit{Working Paper}.

\bibitem[Grantham-McGregor et~al.(2007)]{grantham2007}
Grantham-McGregor, S., Y.~Cheung, S.~Cueto, P.~Glewwe, L.~Richter, and
  B.~Strupp (2007).
\newblock Developmental potential in the first 5 years for children in
  developing countries.
\newblock \textit{Lancet}, 60--70.

\bibitem[Jensen(2010)]{jensen2010}
Jensen, R. (2010).
\newblock The (perceived) returns to education and the demand for schooling.
\newblock \textit{The Quarterly Journal of Economics}, 515--548.

\bibitem[Jopen et~al.(2014)]{jopen2014}
Jopen, G., W.~G\'{o}mez, and H.~Olivera (2014).
\newblock The Peruvian education system: Overview and pending agenda.
\newblock Lima: Working Paper.

\bibitem[Laszlo(2008)]{laszlo2008}
Laszlo, S. (2008).
\newblock Education, labor supply and market development in rural Peru.
\newblock \textit{World Development}, 2421--2439.

\bibitem[Levison and Moe(1998)]{levison1998}
Levison, D., and K.~Moe (1998).
\newblock Household work as a deterrent to schooling: An analysis of adolescent
  girls in Peru.
\newblock \textit{The Journal of Developing Areas}, 339--356.

\bibitem[MINEDU(2019)]{minedu2019}
MINEDU (2019).
\newblock PISA 2018 Assessment.
\newblock Lima.

\bibitem[Miranda(2008)]{miranda2008}
Miranda, L. (2008).
\newblock Factors associated with academic performance and their implications
  for educational policy in Peru.
\newblock \textit{GRADE}, pp.~11--39.

\bibitem[Njie et~al.(2015)]{njie2015}
Njie, H., C.~Manion, and M.~Badjie (2015).
\newblock Girls' family responsibilities and schooling in The Gambia.
\newblock \textit{International Educational Studies}.

\bibitem[Ponce(2012)]{ponce2012}
Ponce, C. (2012).
\newblock Heterogeneous effects of child labor on the acquisition of cognitive
  skills.
\newblock \textit{GRADE}.

\bibitem[Rojas and Cussianovich(2013)]{rojas2013}
Rojas, V., and A.~Cussianovich (2013).
\newblock Growing up in Peru: A longitudinal look at children's time use in
  rural and urban areas.
\newblock \textit{GRADE}.

\bibitem[Romer(1989)]{romer1989}
Romer, P. (1989).
\newblock Human capital and growth: Theory and evidence.
\newblock \textit{NBER Working Papers}.

\bibitem[S\'{a}nchez(2013)]{sanchez2013}
S\'{a}nchez, A. (2013).
\newblock The structural relationship between nutrition, cognitive and
  non-cognitive skills: Evidence from four developing countries.
\newblock \textit{GRADE Institutional}.

\bibitem[Webbink(2005)]{webbink2005}
Webbink, D. (2005).
\newblock Causal effects in education.
\newblock \textit{Journal of Economic Surveys}, 535--560.

\bibitem[Zapata and Contreras(2010)]{zapata2010}
Zapata, D., and D.~Contreras (2010).
\newblock Child labor and schooling in Bolivia: Who's falling behind? The roles
  of domestic work, gender and ethnicity.
\newblock \textit{World Development}, 588--599.

\end{thebibliography}

%=============================================================================
%  APPENDICES
%=============================================================================
\newpage
\appendix

%-----------------------------------------------------------------------------
\section{Young Lives Dataset Structure}
\label{app:data}
%-----------------------------------------------------------------------------

\begin{table}[ht]
\centering
\caption{Young Lives survey rounds: cohorts, ages, and instruments available}
\label{tab:younglives}
\begin{tabular}{cclcccccc}
\toprule
\textbf{Year} & \textbf{Round} & \textbf{Cohort} & \textbf{Age}
  & \multicolumn{2}{c}{\textbf{Reports}}
  & \multicolumn{3}{c}{\textbf{Cognitive tests}} \\
\cmidrule(lr){5-6} \cmidrule(lr){7-9}
 & & & & Child & Caregiver & PPVT & Verbal & Math \\
\midrule
2002 & 1 & Younger & 1  & --- & --- & --- & --- & --- \\
2002 & 1 & Older   & 8  & --- & --- & --- & --- & --- \\
2006 & 2 & Younger & 5  & --- & X   & X   & --- & --- \\
2006 & 2 & Older   & 12 & X   & X   & X   & --- & --- \\
2009 & 3 & Younger & 8  & --- & X   & X   & X   & X   \\
2009 & 3 & Older   & 15 & X   & X   & X   & X   & X   \\
2013 & 4 & Younger & 12 & X   & X   & X   & X   & X   \\
2013 & 4 & Older   & 19 & X   & --- & --- & X   & X   \\
2016 & 5 & Younger & 15 & X   & X   & X   & X   & X   \\
2016 & 5 & Older   & 21 & X   & --- & --- & --- & --- \\
\bottomrule
\multicolumn{9}{l}{\footnotesize Source: Young Lives. Authors' elaboration.}
\end{tabular}
\end{table}

\noindent
\textbf{Note:} Annex~2 of the original thesis contains a map of the survey
sites in Peru (Young Lives, authors' elaboration). [Figure omitted---insert
original figure here.]

%-----------------------------------------------------------------------------
\section{Full Derivation of the Household Optimization Model}
\label{app:model}
%-----------------------------------------------------------------------------

To analyze the relationship between time use and the development of cognitive
skills, I start from the model of time distribution and investment in human
capital \citep{becker1964} and the extension for self-employed households
\citep{gronau1977}. The household is composed of two individuals---caregiver
and child---and has a unitary utility that depends on the consumption of both
individuals in $T$ periods:

\begin{equation}
  \max \; U = u(C_1, C_2, \ldots, C_T)
  \tag{A.1}
  \label{appeq:utility}
\end{equation}

Consumption in each period is composed of the level of cognitive development
of the child ($\NDC_{\child,t}$), domestic tasks ($C_t^{TD}$), and leisure
($C_t^{O}$):

\begin{equation}
  C_t = f_c\!\left(\NDC_{\child,t},\; C_t^{TD},\; C_t^{O}\right)
  \tag{A.2}
  \label{appeq:comp}
\end{equation}

where the leisure and domestic task vectors are defined over both household
members:

\begin{equation}
  C_t^{O} = \begin{pmatrix} c_{\child,t}^{O} \\ c_{\cg,t}^{O} \end{pmatrix},
  \qquad
  C_t^{TD} = \begin{pmatrix} c_{\child,t}^{TD} \\ c_{\cg,t}^{TD} \end{pmatrix}
  \tag{A.3}
  \label{appeq:vectors}
\end{equation}

Following the Beckerian approach, consumption is ``produced'' by household
members through a function $f_j$ that combines the time dedicated to the
activity by both individuals and a set of market goods:

\begin{align}
  C_t^{O}  &= f_{O}\!\left(X_t^{O},\, T_t^{O}\right)
  \tag{A.4}
  \label{appeq:leisure} \\
  C_t^{TD} &= f_{TD}\!\left(X_t^{TD},\, T_t^{TD}\right)
  \tag{A.5}
  \label{appeq:domestic}
\end{align}

where the time vectors are:

\begin{equation}
  T_t^{O} = \begin{pmatrix} T_{\child,t}^{O} \\ T_{\cg,t}^{O} \end{pmatrix},
  \qquad
  T_t^{TD} = \begin{pmatrix} T_{\child,t}^{TD} \\ T_{\cg,t}^{TD} \end{pmatrix}
  \tag{A.6}
  \label{appeq:timevecs}
\end{equation}

Given the emphasis on cognitive development, the level of cognitive development
follows an accumulation process similar to the human capital accumulation
equation of \citet{benporath1967}. The level of cognitive development observed
at period $t$ is equivalent to the level observed at $t-1$ plus the cognitive
development generated at $t$ by the household ($\varphi_t$):

\begin{equation}
  \NDC_{\child,t} = \NDC_{\child,t-1} + \varphi_t
  \tag{A.7}
  \label{appeq:ndc}
\end{equation}

\begin{equation}
  \varphi_t = F(A_0,\, E_t,\, Q,\, N,\, H,\, I)
  \tag{A.8}
  \label{appeq:phi}
\end{equation}

\begin{equation}
  E_t = f_E\!\left(A_0,\; T_t^{E},\; X_t^{E},\; Q,\; N,\; H,\; P_t\right)
  \tag{A.9}
  \label{appeq:education}
\end{equation}

Note that the cognitive development generated in $t$ ($\varphi_t$) depends on
the child's innate ability $A_0$, school performance $E_t$, school quality $Q$,
child characteristics $N$, household characteristics $H$, and the
informational environment $I$. Furthermore, school performance also depends on
time spent in school $T_t^{E}$, a set of goods dedicated to educational
performance $X_t^{E}$, and the price level $P_t$.

On the other hand, market goods used in different activities are obtained from
autonomous income $Y_0$ and the income generated by both individuals through
market work and family work:

\begin{align}
  P_t X_t &= Y_0 + Y_t^{TM} + Y_t^{TF}
  \tag{A.10}
  \label{appeq:budget} \\
  Y_t^{TM} &= f_{TM}\!\left(\hat{W}_t,\, T_t^{TM}\right)
  \tag{A.11}
  \label{appeq:income_tm} \\
  Y_t^{TF} &= f_{TF}\!\left(X_t^{TF},\, T_t^{TF},\, P_t\right)
  \tag{A.12}
  \label{appeq:income_tf}
\end{align}

Finally, a time constraint is incorporated. Total time $T$ is a scarce
resource distributed among all activities by both household members:

\begin{equation}
  T = T_t^{O} + T_t^{TD} + T_t^{E} + T_t^{TF} + T_t^{TM}
  \tag{A.13}
  \label{appeq:time}
\end{equation}

\begin{equation}
  T_t^{j} = \begin{pmatrix} T_{\child,t}^{j} \\ T_{\cg,t}^{j} \end{pmatrix}
  \tag{A.14}
  \label{appeq:timej}
\end{equation}

Solving the household maximization problem, the optimal allocation of time to
each activity depends on exogenous variables:

\begin{equation}
  T_t^{j(*)} = \theta\!\left(\hat{W}_t,\; P_t,\; Q,\; H,\; N,\;
               Y_0,\; A_0\right)
  \qquad \forall\, j \in \{O,\, TD,\, E,\, TF,\, TM\}
  \tag{A.15}
  \label{appeq:optimal}
\end{equation}

\noindent\textbf{Notation:}
\begin{itemize}\setlength{\itemsep}{2pt}
  \item $\NDC_{\child,t}$: cognitive development level (child, period $t$)
  \item $O$: leisure; $TD$: domestic tasks; $E$: education;
        $TM$: market work; $TF$: family work
  \item $A_0$: innate ability (fixed)
  \item $Q$: school quality; $N$: child characteristics; $H$: household
        characteristics; $I$: informational environment
  \item $\hat{W}_t$: market wage; $P_t$: price level; $Y_0$: autonomous income
  \item $(*)$ superscript denotes optimal values
\end{itemize}

\end{document}
